\chapter{Маршрут на два года и развилки}

\section{Календарная модель}
В году 44 учебные недели и восемь резервных. Каждая четвёртая неделя включает
отчёт. Календарные даты заносятся в журнал при начале занятий.
При невыполненном критерии следующий этап переносится.

\section{Первый год: основания и выбор проекта}
\subsection*{Недели 1--12: проверяемая основа}
Материал и задания заданы в предыдущей главе. Особое внимание уделяется
нормировкам, компонентам, компактности и различию двух плоскостей.
\criterion{Вы можете доказать ключевые леммы из квартальной записки у доски
и объяснить, почему две отвергнутые модели репозитория не годятся.}

\subsection*{Недели 13--20: квазиконформная геометрия}
Любич, §§11--18, с.~175--229. Выписать определения и цепочку зависимостей;
полностью разобрать 2--3 доказательства.
Минимум: аналитическое и геометрическое определения qc-отображения,
искажение модуля, нормировка решения уравнения Бельтрами, смысл голоморфного
движения и его продолжения. MRMT и теорему продолжения можно сначала
использовать как явно названные внешние теоремы. Полное доказательство MRMT
выделить в отдельный блок, если оно нужно выбранной теме.

Вычисления: проверить радиальное растяжение \(z\mapsto z|z|^{\alpha-1}\)
при \(\alpha>0\), его дилатацию и действие на модуль кольца.
\criterion{Письменный расчёт qc-примера, точная теорема об искажении модуля
и одна схема склейки с обозначенными входными условиями.}

\subsection*{Недели 21--28: динамические основания и первый отбор тем}
Любич, §§19--24, с.~234--344, выборочно. Нужны классификация периодических
точек, нормальные семейства, заполненное множество Жюли, координата Бёттхера
у бесконечности и роль посадки рациональных лучей. На этом этапе различать
локальное построение координаты и её глобальное продолжение, а также
утверждения для отдельных лучей и для всей границы.

К концу 24-й недели составить два исследовательских паспорта и сопоставить
их с существующими декларациями проекта.
\criterion{Один полный аргумент из динамики и два точных кандидата на
исследовательский вопрос с ближайшими источниками.}

\subsection*{Недели 29--36: puzzle, polynomial-like и параметрический переход}
Обзор определений из §§25, 28, 31, 33--34, 38, 40, 42, 47 книги Любича.
Для полного разбора выбрать один механизм: компактное
сжатие корректных параметрических кусочков, конкретный конечный puzzle
или локальную polynomial-like конструкцию.

Нарисовать отдельно фазовый и параметрический графы. Подписать, что именно
движется с параметром, где использовано голоморфное движение и откуда
берётся связность параметрической части. При невозможности доказать шаг
указать точную литературу и гипотезы. Сравнить с \repo{ParameterPuzzleBasis}
из текущих \repo{Mlc/Literature/} модулей.
\criterion{Доклад 30--40 минут, один проверенный механизм и выбор единственной
основной ветки A, B или C после внешнего обсуждения.}

\subsection*{Недели 37--44: восьминедельный исследовательский пилот}
У пилота одна центральная лемма и один класс примеров. Каждые две недели
должен появляться проверяемый объект: доказательство, опровергнутая гипотеза,
корректный сертификат или уменьшенная постановка. К окончанию пилота написать
результат как короткую статью: постановка, связь с известным, определения,
основное утверждение, доказательство, пример и ограничения.
\criterion{Законченная записка по пилоту и задача на следующий цикл.}

\section{Второй год: один вклад и законченная рукопись}
\begin{table}[!htp]
\centering\small
\caption{Второй год: этапы и результаты}
\begin{tabularx}{\textwidth}{@{}P{0.13\textwidth}P{0.4\textwidth}Y@{}}
\toprule
Недели & Работа & Условие перехода\\
\midrule
45--56 & Уточнение основной конструкции; проверка новизны; исправление
семантики и один нетривиальный пример & Зафиксированы утверждение, точный
класс объектов, ближайший результат и согласованный внешний отзыв\\
57--68 & Главная лемма; оценки и контрпримеры к лишним обобщениям &
Есть полный письменный аргумент либо доказана причина отказа от него;
следующий цикл имеет изменённую постановку\\
69--76 & Независимая проверка; выбранная Lean-реализация; воспроизводимые
сертификаты & Замечания разобраны, проверенные части отделены от условных\\
77--84 & Сборка рукописи; введение переписывается под фактический результат &
В тексте нет неявных переходов между плоскостями, неподтверждённой новизны
и пустых гипотез под именем содержательной теоремы\\
85--88 & Два пробных доклада; ответы на возражения; решение о подаче или
третьем годе & Внешний математик понимает вклад, вы можете защитить ключевое
доказательство без ИИ; дальнейшие требования согласованы отдельно\\
\bottomrule
\end{tabularx}
\end{table}

На неделе~68 по результатам рецензии выбрать продолжение: исходная задача,
более узкая постановка или дополнительный исследовательский год.

\section{Алгебраическая развилка B}
Если эта ветка оказывается ближе вашим интересам, после общего квартала
заменить часть аналитического маршрута на 12--16 недель следующей цепочки:
\begin{enumerate}
\item Накрытия и подъём путей; фундаментальная группа проколотой сферы;
      действие монодромии. Самостоятельно вычислить случай \(z\mapsto z^2\).
\item Коммутирующие левое и правое действия; тензорное произведение бисетов;
      зависимость от базовой точки и выбранных путей.
\item Bartholdi--Dudko I~\cite{BDI}: определения и один ограниченный случай
      теоремы ван Кампена для бисетов.
\item Один пример конечного посткритического типа, например \(z^2-1\),
      с полностью описанными маркировками и действиями.
\item Затем только нужный фрагмент IV/V~\cite{BDIV}, если исследовательский
      вопрос действительно касается расширяющей структуры.
\end{enumerate}
Групповые ряды и гомологические инструменты изучать по зависимостям
выбранного утверждения.

\section{Развилка C: современная аналитическая геометрия}
После освоения основных инструментов выбрать один маршрут:
\begin{enumerate}
\item конечная ренормализуемость и Yoccoz;
\item quasi-additivity и одна лемма переноса модуля Kahn--Lyubich;
\item одна указанная гипотеза теоремы о спутниковых или Feigenbaum-параметрах.
\end{enumerate}
Сначала восстановить ограниченную известную цепочку, затем искать усиление.
По обзору Дудко~\cite{DudkoSurvey} составить карту зависимостей.
Работы о спутниковых параметрах и Feigenbaum-точках
дают отдельные классы локальной связности~\cite{Satellite,Feigenbaum}.
Работа Kahn--Kapiamba--Lyubich 2026 года относится к параболически ограниченной
примитивной ренормализации~\cite{KKL}.
